\documentclass[../Main.tex]{subfiles}

\begin{document}

\section{Evaluation} \label{sec:Evaluation}

Das Ziel des Playtests ist, die acht Designvorschläge (V1--V8), die wir aus unseren Spieleanalysen und der Leveldesign-Recherche hergeleitet und in Kapitel~\ref{sec:Umsetzung} umgesetzt haben, zu überprüfen, also zu prüfen, ob diese die beabsichtigten Wirkungen auf Spielfluss, Schwierigkeit und Spielgefühl tatsächlich erzeugen.

Wir sind beim Playtest wie in Abschnitt~\ref{sec:VorgehenPlaytest} beschrieben vorgegangen: Zuerst wurde die Testperson mit dem Pre-Questionnaire eingeschätzt, danach folgte der Playtest, der aufgenommen wurde, und anschliessend das Post-Questionnaire sowie das semistrukturierte Interview für tiefergehende Fragen.

Das folgende Kapitel beschreibt zunächst den Ablauf der Playtests und geht dann auf die für die Erkenntnisfrage relevanten Ergebnisse ein. Darauf folgt eine Diskussion, in der wir die Designvorschläge einzeln durchgehen, und schliesslich unsere Empfehlungen für weitere Playtests.

\begin{itemize}[leftmargin=1.5cm]
    \item \textcolor{red}{Abschnitt~\ref{sec:Verbesserungen} heisst \enquote{Verbesserungsvorschläge}. Meint ihr Empfehlungen für weitere Playtests oder Verbesserungen am Spiel? Begriff angleichen.}
\end{itemize}

\subsection{Durchführung und Testpersonen} \label{sec:Durchfuehrung}

Die Playtests wurden immer von Rodin Arslan durchgeführt; Zhile Zhu war hin und wieder dabei. Gespielt wurde immer mit Maus und Tastatur. Drei Playtests fanden über Discord mit Bildschirmübertragung statt, zwei vor Ort. Alle Testpersonen spielten dieselbe exportierte Version des Spiels (für Mac oder Windows) und benutzten dabei ihre eigene, gewohnte Hardware.

Während des Spielens wurde kein Smalltalk geführt; in entspannten Phasen haben wir mit den Testpersonen allerdings über das Spiel gesprochen. Hatte eine Testperson eine Frage zum Kampfsystem oder etwas vergessen, wurde ihre Frage beantwortet. Ganz zu Beginn wurde den Testpersonen zudem gezeigt, welche Aktionen es gibt und wie man sie ausführt, da der Prototyp kein Tutorial enthält.

\begin{itemize}[leftmargin=1.5cm]
    \item \textcolor{red}{Ergänzen: Zeitraum der Playtests, ungefähre Dauer pro Testperson, wie aufgenommen wurde (auch bei den Discord-Tests).}
    \item \textcolor{red}{Abweichungen vom Plan kurz benennen: Einführung durch uns statt Tutorial, Gespräche in entspannten Phasen, Remote-Tests über Discord (Beobachtung nur über den Bildschirm). Diese Punkte auch in den Limitationen (Kapitel~\ref{sec:Schlussfolgerungen}) aufgreifen.}
    \item \textcolor{red}{Anonymisierte Tabelle der Testpersonen TP1--TP5: Spielzeit pro Woche (P1), Lieblingsgenres (P2), Selbsteinschätzung (P5), Spielstil (P6); dazu 2--3 Sätze zur Zusammensetzung.}
    \item \textcolor{red}{Die alte Ablaufbeschreibung (Pre-Questionnaire aushändigen usw.) habe ich gelöscht, sie steht bereits in Abschnitt~\ref{sec:MethodePlaytest}, inklusive der Abweichung vom Think-aloud.}
\end{itemize}

\subsection{Ergebnisse}

Die einzelnen Antworten der Testpersonen sind im Anhang aufgeführt (Abschnitt~\ref{sec:PlaytestProtokolle}). Im Folgenden werden die Ergebnisse pro Designvorschlag beschrieben.

\subsubsection{Ergebnisse pro Designvorschlag} \label{sec:ErgVorschlag}

\paragraph{V1: Überlegenheit als Grundzustand}
Der Aussage, sich den Gegnern grundsätzlich überlegen gefühlt zu haben (Q1), stimmten vor allem jene Testpersonen zu, die auch sehr aggressiv gespielt haben. TP3 stimmte beispielsweise völlig zu; diese Testperson ignorierte auch ihre eigenen Lebenspunkte und konzentrierte sich nur auf die Gegner. TP5 spielte dagegen weniger aggressiv, eher taktisch und methodisch, und bewertete die Aussage neutral. TP4 meinte: \enquote{Ich geniesse meine Überlegenheit}; besonders befriedigend sei es, alle Gegner auf einmal zu besiegen. TP5 hätte sich den Spieler im Vergleich zu den Gegnern hingegen noch stärker vorgestellt, etwa wie in \game{Hyrule Warriors}. Insgesamt gaben drei von vier Testpersonen an, sich überlegen gefühlt zu haben.

Die Überlegenheit wurde aber auch immer wieder durchbrochen. Insbesondere wenn die Testpersonen von einer Gegnergruppe umzingelt waren, konnte es passieren, dass sie wie in einem Stunlock gefangen waren und erst nach erlittenem Schaden wieder herauskamen. In der Luft waren sie zudem durch die Archers verwundbar: Wurden sie getroffen, fielen sie zu Boden, und es kam vor, dass sie beim Aufstehen weiteren Schaden erlitten. Auch die Trägheit der Combos machte die Testpersonen anfällig. Während einer Combo spielte dies eine geringere Rolle, weil die Gegner den Spieler dann kaum angreifen konnten; am Ende einer Combo konnte es jedoch passieren, dass man Schaden erlitt. Kam während einer Combo ein Pfeil angeflogen, konnte man fast nicht ausweichen, und viele Testpersonen haben den Schaden einfach in Kauf genommen.

\begin{itemize}[leftmargin=1.5cm]
    \item \textcolor{red}{Verweise auf den Anhang ergänzen, z.\,B. TP3 (B1), TP4 (B11), TP5 (Q14), Stunlock (TP5 Q13, TP4 B2), Treffer in der Luft (TP2 I6, TP4 B11, TP5 I8), Combos (TP3 I6). Gilt für alle Vorschläge.}
    \item \textcolor{red}{\enquote{Während einer Combo \ldots{} weil die Gegner den Spieler dann kaum angreifen konnten}: Euer Satz war hier unklar (\enquote{weil nicht einen angreifen konnte}). Prüfen, ob es so gemeint ist.}
    \item \textcolor{red}{\enquote{viele Testpersonen} $\rightarrow$ konkret: drei von vier tankten Treffer (TP2, TP3, TP4; B4).}
    \item \textcolor{red}{Fehlt: Der Heavy Soldier, der die Überlegenheit gezielt brechen soll, wurde als eher harmlos erlebt (TP4 I2: verfehlt die Hälfte der Zeit; TP3 I5 und B3; TP2 B1 tankt ihn). Wichtig für die Diskussion.}
\end{itemize}

\paragraph{V2: Begrenzung gleichzeitiger Angreifer}
Alle vier Testpersonen stimmten der Aussage, auch bei vielen Gegnern die Kontrolle behalten zu haben (Q2), eher oder völlig zu. Sie hatten dabei alle eine eigene Strategie, um mit grösseren Gegnermassen umzugehen. Einerseits liessen sich ähnliche Muster beobachten: Einige griffen einfach den ersten Gegner an, der bei ihnen war, andere wollten zuerst alle Archers loswerden. Andererseits gab es unterschiedliche Ausprägungen: TP2 rannte weg, um sich heilen zu können, während TP5 auf die Boxen sprang, die er gleich zu Beginn eines Encounters von Gegnern befreit hatte, um einen sicheren Rückzugsort zu haben. TP3 kämpfte immer gegen viele Gegner gleichzeitig; dabei kam es vor, dass die Testperson einen kurzen Moment der Panik hatte, wenn sie zu wenig Leben hatte. TP5 versuchte eher, Gruppen zu meiden und nur einzelne Gegner anzugreifen.

TP2 meinte, es sei schwerer, wenn die Gegner in einem Haufen seien, als wenn sie einzeln kämen, aber nicht unangenehm schwer. TP5 meinte, er könne nicht einschätzen, wie viel Schaden eine Gruppe von Gegnern mache, wenn er in sie hineinlaufe, und ob er dann im Haufen auch noch in einen Stunlock gerate \textcolor{red}{[Satz endete mit \enquote{und deswegen} -- vervollständigen]}.

\begin{itemize}[leftmargin=1.5cm]
    \item \textcolor{red}{Korrigiert: In eurem Text stand zweimal \enquote{TP2}. Das Säubern der Boxen als Rückzugsort war laut Anhang TP5 (B3), bitte prüfen.}
    \item \textcolor{red}{Fehlt: TP4 häufte Soldiers beim Archer an, um alle auf einmal zu treffen (B2).}
    \item \textcolor{red}{Der Satz zu den Mustern (erster Gegner / zuerst Archers) gehört eher zu V3; hier ggf. kürzen.}
\end{itemize}

\paragraph{V3: Ein taktisches Dilemma durch zwei Gegnerrollen}
Im Interview gaben alle Testpersonen an, zuerst die Archers anzugreifen, wenn sich die Möglichkeit bietet. TP2 meinte, sie sei vor allem auf die Chain und Heavy Archers losgegangen und habe sich die Heavy Soldiers für den Schluss aufbewahrt, da die Archers mit höherer Wahrscheinlichkeit trafen und die anderen Gegner besser kontrollierbar waren. TP3 meinte, sie müsse nicht unbedingt auf die Soldiers losgehen, da diese einen ohnehin suchen und auf den Spieler zurennen. TP4 hat das taktische Dilemma zwar gesehen, meinte aber, es habe für ihn keine allzu grosse Rolle gespielt; er habe einfach, wie er wollte, alles kurz und klein gehauen. Erst wenn er merkte, dass es so nicht ging, und dazu gezwungen war, habe er sich Gedanken zum Vorgehen gemacht.

Beobachtet wurde vor allem bei TP2 und TP3 das erwartete Verhalten: zuerst die Archers, dann die Soldiers, wobei die Soldiers den Weg zu den Archers blockierten. TP4 schlug dagegen zuerst auf die Soldiers ein und ging erst danach auf die Archers los. TP5 probierte Verschiedenes aus: Manchmal griff er zuerst die Archers an, manchmal die Nahkämpfer. Ausschlaggebend waren für ihn eher die Gegner auf den Boxen, weil er sich so einen sicheren Platz zum Heilen verschaffen konnte.

\begin{itemize}[leftmargin=1.5cm]
    \item \textcolor{red}{Angepasst: Euer erster Satz lautete, das Dilemma sei von allen \emph{wahrgenommen und beobachtet} worden, dass alle zuerst auf die Archers gehen. Im Interview sagen das alle (I5), beobachtet wurde es aber nur bei TP2 und TP3 klar. Q3 als Tendenz ergänzen (TP4 neutral, die anderen stimmen zu).}
    \item \textcolor{red}{TP2 laut Anhang: zuerst Archers mit drei Pfeilen oder zielsuchenden Pfeilen; die \enquote{mit den Hämmern} zum Schluss, weil langsam und konterbar (I5). \enquote{Ignorierte} habe ich daher zu \enquote{für den Schluss aufbewahrt} geändert.}
\end{itemize}

\paragraph{V4: Klar telegrafierte Angriffe}
Bei der Aussage, die Angriffe seien klar angekündigt worden (Q4), gehen die Meinungen am weitesten auseinander. Interessanterweise stimmte TP3, die sehr aggressiv und ohne viel nachzudenken spielte, völlig zu. TP2, die ebenfalls sehr aggressiv spielte und meistens erst reagierte, nachdem ein Treffer gelandet war, bewertete die Aussage neutral. Auch TP5, der eher taktisch spielte, bewertete sie neutral, da er die Ankündigungsanimationen der Archers verwirrend fand. TP4, der Projektilen auswich, wenn er konnte, sonst aber alle Treffer in Kauf nahm, stimmte überhaupt nicht zu. Dazu meinte er, die Archers seien besonders, da sie ausserhalb des Blickfelds angreifen könnten und es keinen Sound gebe, der den Angriff ankündigt. TP5 meinte zudem, die Tode seien wegen der Ankündigungen nicht unfair gewesen: Man habe kommen sehen, dass man stirbt, und es habe sich eher wie ein selbstverschuldetes Problem angefühlt.

\begin{itemize}[leftmargin=1.5cm]
    \item \textcolor{red}{Verschoben: Eure Erklärungen (\enquote{das kann damit zusammenhängen \ldots}, \enquote{das könnte darauf hindeuten \ldots}) sind Interpretationen und stehen jetzt in Abschnitt~\ref{sec:Interpretation} unter V4.}
    \item \textcolor{red}{\enquote{ausserhalb des Blickfelds}: Im Anhang steht bei TP4 nur \enquote{ohne wirkliche Prediction oder Sound-Effekte} (I5). Falls das aus der Aufnahme stammt, im Anhang ergänzen.}
    \item \textcolor{red}{Ergänzen: TP3 sah die Pfeile kommen, konnte wegen der Combos aber nicht ausweichen (I6, Q13).}
\end{itemize}

\paragraph{V5: Lesbare Platzierung, dynamisches Verhalten}
Hier stimmten TP2 und TP3 eher zu, TP4 und TP5 eher nicht (Q5). TP2 kletterte zwischen den Wellen auf Boxen, um den Raum zu überblicken. TP3 griff einfach an, was vor ihr stand; die vorhandene Information reichte ihr also aus. In Raum~10, einer engen, verwinkelten Gasse, fand TP3 als einzige Testperson den versteckten Archer eine Zeit lang nicht. Auch TP5 griff eher an, was vor ihm lag, und verstand das dynamische Verhalten der Gegner. Er meinte, man wisse nicht, wo die Gegner spawnen, bevor man den Raum betreten habe. TP5 kreiste zudem hin und wieder eine Weile im Raum, bevor er angriff, um die Gegner zu beobachten, da er so nicht getroffen werden konnte; auf diese Weise sammelte er die Gegner auch an. TP4 fand es unangenehm, von hinten, also ausserhalb des Sichtfelds, getroffen zu werden. Er schaute sich eine Welle eher lange an, bevor er angriff.

TP3 bemerkte, dass Gegner ausserhalb des Sichtfelds spawnen, meinte aber, man bemerke sie früh genug, um reagieren zu können. Sie stellte auch fest, dass die Archers eher erhöht spawnen und dass es definitiv eine manuell platzierte Konstante geben müsse; es seien also deutliche Muster erkennbar. TP2 gefiel die Verteilung der Gegner im Raum am besten. TP5 meinte nach einem Tod, es fühle sich richtig gut an, wenn man wisse, wo ein Gegner spawnen wird, und ihn direkt zu Beginn ausschalten könne.

\begin{itemize}[leftmargin=1.5cm]
    \item \textcolor{red}{Ergänzen: TP4 starb in Raum~9, weil er einen Archer übersehen hatte (B6). Optional: TP5 wurde von der Gegnermasse fast in einen Türrahmen gedrängt (B6).}
    \item \textcolor{red}{Verweise auf den Anhang ergänzen: TP2 (B6, Q12), TP3 (B6, I8), TP4 (B6, B11), TP5 (B2, B6, Q12).}
    \item \textcolor{red}{\enquote{Raum~10}: Im Anhang steht bei TP3 nur \enquote{versteckter Heavy Archer} ohne Raumnummer. Falls aus der Aufnahme, dort ergänzen.}
\end{itemize}

\paragraph{V6: Der Raum als Druckmittel}
Hier gehen die Meinungen etwas auseinander. TP2 spielte in grösseren und kleineren Räumen unterschiedlich: In grossen Räumen spielte sie aggressiver, in kleineren weniger aggressiv und eher taktisch, da es dort wenig Platz zum Heilen gab und das Leben dadurch wertvoller war. Bei TP3 war es unterschiedlich: In manchen kleinen Räumen spielte sie genau gleich aggressiv, in anderen hektischer. TP5 bemerkte dagegen einen deutlichen Unterschied in kleinen Räumen. Wenn diese verwinkelt waren, empfand er sie als gruselig und schaute deshalb vorsichtig um die Ecken.

TP2 meinte, es gebe eine gute Mischung der Raumgrössen. Die kleinen Räume seien tendenziell schwerer, weil es keine Zufluchtsmöglichkeit gebe; ihr Lieblingsraum sei aber Raum~6 mit der Brücke gewesen. TP3 meinte, kleine Räume seien einfacher zu überblicken gewesen, man habe direkt gewusst, was gespawnt sei und was auf einen zukomme. In grösseren Räumen mit Boxen und Brücke habe es dagegen mehr Möglichkeiten gegeben, sich im Raum zu bewegen. Für TP4 spielte die Raumgrösse keine Rolle, weil er einfach alles kopfvoran versucht hat. TP5 meinte, kleine Räume seien übersichtlicher, man könne aber in Ecken gedrängt werden; die grossen seien schwerer gewesen, weil es dort mehr Gegner gab. Wie bereits bei V2 beschrieben, hatte TP5 etwas Mühe mit grösseren Gegnermengen.

\begin{itemize}[leftmargin=1.5cm]
    \item \textcolor{red}{Korrigiert: \enquote{für TP2 spielte es keine Rolle} war laut Anhang TP4 (I4), bitte prüfen.}
    \item \textcolor{red}{Verschoben: Eure Erklärung zu TP3 (\enquote{das liegt wahrscheinlich daran \ldots}) steht jetzt in Abschnitt~\ref{sec:Interpretation} unter V6.}
    \item \textcolor{red}{Q6 als Tendenz ergänzen: Die Antworten gehen weit auseinander (TP2 und TP5 stimmen völlig zu, TP3 neutral, TP4 eher nicht).}
    \item \textcolor{red}{Raum~6 mit der Brücke: Im Anhang nennt TP2 \enquote{Raum 7.5}. Einheitlich machen. \enquote{gruselig} (TP5) steht nicht im Anhang -- falls aus der Aufnahme, dort ergänzen.}
\end{itemize}

\paragraph{V7: Wellenstruktur mit Erholungsphasen}
Alle vier Testpersonen organisierten ihre Erholungszeit selbst. TP2 machte Pause, bevor sie in den nächsten Raum ging, und nutzte die Boxen. TP3 ist gestorben, weil sie sich anfangs nicht geheilt hat; nachdem sie das bemerkt hatte, legte sie am Ende eines Raumes eine Ruhephase zum Heilen ein. Auch TP4 heilte sich vor dem Raumwechsel. TP5 ging dabei sogar etwas taktischer vor und liess oft einen Gegner am Ende einer Welle am Leben, um sich vor der nächsten Welle zu heilen, da man nicht genau wusste, was danach kommt.

TP2 meinte, es sei schwer, während eines Encounters Heilpausen zu finden. TP3 meinte, es gebe eine gute Balance; sobald alle Fernkämpfer besiegt seien, werde es entspannt. TP4 empfand den Verlauf der Kampfabschnitte ebenfalls als angenehm: Er sei immer auf das Spiel konzentriert gewesen, und nichts habe seine Vertiefung gestört. TP5 meinte, die Spannung hänge vom eigenen Leben ab, da man mit wenig Leben angespannter sei.

\begin{itemize}[leftmargin=1.5cm]
    \item \textcolor{red}{Verschoben: \enquote{Das liegt daran, dass TP2 die Wellenstruktur nicht erkannt hat} ist eine Interpretation und steht jetzt in Abschnitt~\ref{sec:Interpretation} unter V7.}
    \item \textcolor{red}{TP5 \enquote{liess oft einen Gegner am Leben}: Im Anhang steht \enquote{heilt sich, wenn am Ende einer Welle nur noch wenige Gegner übrig sind} (B8). Prüfen und angleichen.}
    \item \textcolor{red}{Q7 als Tendenz ergänzen (TP3 und TP5 stimmen völlig zu, TP2 und TP4 neutral). TP2 sprach zudem von \enquote{überlappenden Wellen} (I7) -- wichtig für die Diskussion.}
\end{itemize}

\paragraph{V8: Gestaffelte Einführung neuer Gegner}
Hier stimmten alle Testpersonen völlig zu (Q8). Alle erkannten neue Gegnertypen bei ihrer Einführung an der Farbe, am Aussehen oder an ihren Angriffen. Wie sie mit ihnen umgingen, war allerdings unterschiedlich: TP2 hielt meist kurz inne und analysierte, ob sich noch etwas im Raum befand. TP4 beobachtete den Chain Archer zuerst, ohne anzugreifen, um zu sehen, was er macht. TP3 und TP5 griffen neue Gegner sofort an. TP3 dachte bei der Einführung des Archers, es handle sich einfach um eine stärkere Version der Soldiers, und wusste anhand des Aussehens anfangs nicht, was die neuen Gegner tun.

\begin{itemize}[leftmargin=1.5cm]
    \item \textcolor{red}{Allgemein für 6.2.1: Pronomen einheitlich halten (ich habe für TP2 und TP3 \enquote{sie} wie im Anhang verwendet, für TP4 und TP5 \enquote{er}). Für die Anonymisierung evtl.\ durchgehend \enquote{die Testperson}.}
    \item \textcolor{red}{Sobald TP1 da ist: \enquote{x von 4} zu \enquote{x von 5} anpassen.}
\end{itemize}

\subsubsection{Gesamteindruck und allgemeine Beobachtungen} \label{sec:ErgAllgemein}

Der Gesamteindruck der Testpersonen in den drei Kategorien fällt ziemlich ähnlich aus. Das Spielgefühl beurteilten alle als eher befriedigend (Q9). Den Spielfluss empfanden ebenfalls alle als zusammenhängend (Q10); dabei meinten sie, dass die Räume zwar jeweils für sich stehen und wenig miteinander verbindet, der Spielfluss insgesamt aber gepasst habe. Auch die Schwierigkeit fanden alle angemessen (Q11): Alle sind ein- bis zweimal gestorben, aber nicht wiederholt an derselben Stelle, und es gab viele Momente, in denen sie knapp überlebten.

Als das, was sich am besten anfühlte (Q12), nannten die Testpersonen die Gegnerverteilung und die Spawns, etwa wenn man Gegner direkt beim Spawnen besiegen kann. Auch die taktischen Entscheidungen, das Heilen und das Weglaufen wurden als interessant genannt. Horden zu vernichten, was der Kern des Spiels ist, habe Spass gemacht und sich gut angefühlt. Positiv erwähnt wurde zudem die Tiefe durch die unterschiedlichen Combos, also Light und Heavy mit ihren verschiedenen Schadenswerten.

Als unfair oder unangenehm empfanden die Testpersonen (Q13), dass man Combos nicht abbrechen und deshalb nicht so einfach auf Gegner reagieren konnte. Auch in einer Horde von Gegnern eingeklemmt zu sein, fühlte sich unbefriedigend an. Eine Testperson meinte zudem, das Springen fühle sich nicht wie erwartet an; das deckt sich mit einer anderen Aussage, wonach der Übergang zwischen Bewegung und Kampf etwas unstimmig sei.

Im Vergleich mit ihren Erwartungen (Q14) waren die Testpersonen vom Umfang und von der Gegnervielfalt positiv überrascht. Erwartet wurde ein Boss, der nicht kam; zudem wurden hin und wieder Moves wie Block oder Dodge gewünscht. Es gab aber auch Stimmen, die ihre Meinung am Ende änderten und meinten, dass diese Moves nicht ins Spiel gepasst hätten. Ausserdem wurde erwartet, dass man als Spieler noch stärker sei.

Weitere Beobachtungen betreffen die Orientierung, die nicht immer reibungslos verlief: Oft wurde nicht bemerkt, dass ein Encounter beendet war, und offene Türen wurden übersehen, da es keinen audiovisuellen Hinweis darauf gab. Bei der Steuerung war den Testpersonen die Tiefe der Eingaben nicht immer intuitiv. So kann man die Dash-Taste gedrückt halten; Testpersonen, die das aus anderen Spielen kannten, waren davon begeistert, während die anderen die Taste eher spammten. Beim Springen verhielt es sich ähnlich. Auch die Reichweite der eigenen Waffe einzuschätzen, fiel etwas schwer, und die Dauer der Heilung war nicht ganz nachvollziehbar. TP2 schlug vor, statt der mündlichen Erklärung ein Tutorial einzubauen. TP4 und TP5 versuchten, mit Objekten zu interagieren, und suchten beide nach Glitches.

Die Stimmung war in allen Playtests positiv: TP3 und TP4 lachten viel, und TP3 und TP5 spielten nach dem Ende einfach weiter. TP2 war hin und wieder genervt, wenn sie beim Heilen unterbrochen wurde. Fast jede Testperson hatte einen Gegnertyp, der sie besonders störte. TP5 meinte, er sei nie unter Stress gewesen.

\begin{itemize}[leftmargin=1.5cm]
    \item \textcolor{red}{\textbf{Faktencheck mit dem Anhang:}}
    \item \textcolor{red}{\enquote{Alle sind ein- bis zweimal gestorben}: TP2 ist laut Interview gar nicht gestorben (I3), TP5 dreimal (I3), TP4 mindestens zweimal (I1). Anpassen.}
    \item \textcolor{red}{\enquote{Den Spielfluss empfanden alle als zusammenhängend}: TP3 gab \enquote{1 von Raum zu Raum, 4 im Kampf} an (Q10). Das deckt sich mit eurem Satz über die getrennten Räume -- so formulieren, statt \enquote{alle}.}
    \item \textcolor{red}{Glitches suchten laut Anhang TP3 (I7: \enquote{Bugs zu finden hat mir Freude bereitet}) und TP5 (B13), nicht TP4. TP4 und TP5 wollten mit Objekten interagieren (B13).}
    \item \textcolor{red}{Nicht im Anhang belegt: TP4 lachte viel, TP5 spielte nach dem Ende weiter, Testpersonen änderten ihre Meinung zu Block/Dodge, Begeisterung über das Gedrückthalten von Dash. Falls aus den Aufnahmen, im Anhang ergänzen.}
    \item \textcolor{red}{Pauschale Aussagen (\enquote{die Testpersonen}, \enquote{oft}) durch \enquote{x von 4} bzw.\ TP-Nummern mit Verweis ersetzen, z.\,B. offene Türen übersehen: TP2 und TP4 (B10), TP5 nicht (B14); Springen: TP2 (Q13); Übergang Bewegung/Kampf: TP5 (B10); Reichweite: TP3 (Q13), TP5 (B10); Heilungsdauer: TP4 (B10); Boss erwartet: TP2 (Q14), TP3 (I2).}
\end{itemize}

\subsection{Diskussion}

\subsubsection{Interpretation pro Designvorschlag} \label{sec:Interpretation}

\begin{itemize}[leftmargin=1.5cm]
    \item \textcolor{red}{Pro V1--V8: bestätigt / teilweise bestätigt / nicht bestätigt, mit Begründung aus den Daten und Bezug zur Recherche. Abweichungen zwischen Vorschlag und Umsetzung (z.\,B. V2, V7) hier diskutieren. Selbst schreiben.}
    \item \textcolor{red}{Unten stehen eure Erklärungen, die ich aus 6.2.1 hierher verschoben habe. Sie sind ein guter Anfang für V4, V6 und V7.}
\end{itemize}
\paragraph{V1}
Bei den aggressiven spielern trifft das völlig zu. sie geniessen ihre überlegenheit. dagegen ist bei den eher taktischen spielern der spieler noch etwas zu schwach und könnte stärker sein. die einzelnen gegner sind an sich nie ein problem wie geplant nur in combinaiton und in grossen mengen.erwartet war, dass der heavy Soldier ein bruch in der überlegenheit ist, wir haben ihn auch gut vorgestellt in einem kleinen raum aber im verlauf der playtests haben alle tester gemerkt, dass er wenig akkurat mit seinen treffern ist und deshalb war er weniger eine bedrohung. die überlegenheit ist aber durch weitere ungeplante dinge gebrochen worden. die trägheit der kombos und auch treffer in der luft waren solche auch waren die Gegnergruppe, da wir sie sehr agressiv eingestellt haben auch eine grössere herausforderung als erwartet. im grundzustand konnten wir diese these bestätigen, die spieler spielten agressiv.das gezielte durchbrechen lief nicht ganz wie geplant, da der heavy Soldier etwas zu schwach geraten ist. Dass wir die Überlegenheit durch die anderen wege brechen hatten wir schon vor dem playtest bemerkt aber haben diese eher schwächer und nur ergänzent eingeschätzt als dass sie die tragenden brecher sind. Die testpersonen haben den Kontrollzustand wie bei Nier Automata verstanden also als völlige kontrolle als Grundzustand. Des weiteren gab es wie 2 bewegungs modi einmal die normale vortbewegung und dann noch kämpfen und der switch zwischen diesen modi hat den spielern hin und wieder schwierigkeiten bereitet und die reichweite des spielers war etwas schwer einzuschätzen.

\paragraph{V2}
Dem stimmen die Testpersonen alle eher oder völlig zu. es wurde als schwer aber nicht unangehnehm beschrieben. Unsere umsetzung von grossen und agressiven gruppen hat aber auch dazu geführt, dass man nicht so genau wusste, wie viel schaden so eine gruppe abziehen kann. wenn man in der luft getroffen wird und dann wieder aufstehen muss, kann es sein, dass man von einer gruppe umzingelt wird und viel schaden erleidet, was vielleicht zu sehr punishing ist. Wir haben in unserer umsetzung die anzahl der gespawnten gegner limitiert und weniger wie viele von denen aktiv sind und schlagen. das hat sich vor dem playtest besser angefühlt. es könnte aber auch sein, dass wir damit übertireben haben. Mit diesem vorschlag wollten wir erreichen, dass sich der combat schwer anfühlt aber die schwierigkeit dennoch machbar ist. schlussendlich haben wir es etwa wie Ninja Gaiden gelöst, wo die maximale anzahl gegner begrenzt wird.

\paragraph{V3}
das taktische Dilemma wurde von allen spieltestern wahrgenommen und es gab verschiedene ansätze, dieses Dilemma zu lösen. grundsetzlich konnte man beobachten, das Archers und die Boxen zu räumen vorrang hatten. Aufgrund der unterschiedlichen typen von gegnern hatten die individuellen playtester allerings unterschiedliche gegner, die sie als lästig empfunden haben. Wir konnten allerdings auch beobachten, dass man das spiel durchspielen konnte mit nur wenig taktik. Es gab auch wiedersprüchliche antworten von den spielern, manche haben zwar das taktische Dilemma gesehen aber nicht immer demnach gehandelt. so war die herangehensweise einfach den erst besten gegner zu schlagen oft auch zielführend und schaden von archern wurde teils ignoriert. Dieser Vorschlag wurde ziemlich gut umgesetzt da archers oft händisch platziert worden sind und oft auch am rand. da sie sich nicht bewegt haben war so oft ein Soldier zwischen spieler und archer. Ob dann der archer zuerst anvisiert wurde oder nicht war dann halt untershciedlich oder das dilemma wurde erkannt aber in manchen situationen war schaden nehmen und durch die mental einfachere option und passte auch zum agressiven spielstyle.

\paragraph{V4: Klar telegrafierte Angriffe}
Dass TP3 völlig zustimmte, kann damit zusammenhängen, dass ihr die meisten Treffer egal waren und ihr die Ankündigungen deshalb genügten. Dass TP2 erst nach einem Treffer reagierte und nur neutral bewertete, könnte darauf hindeuten, dass die Ankündigungszeit etwas zu kurz war. Dass TP4 überhaupt nicht zustimmte, kann damit zusammenhängen, dass er die Information nicht brauchte und deshalb nicht darauf achtete, oder auch daran liegen, dass die Combos zu träge sind und er deshalb auf die Ankündigungen nicht reagieren konnte. Tp5 meinte, dass das telegraphing und dass man dadurch vorhersehen konnte dass man stirbt, den tot nicht unfair sondern weniger schlim gemacht haben, da dieser nicht aus dem nichts kam. Hier gab es allerdings die grösste streuung in dne antworten der teilnehmer. das liegt wahrscheinlich an ihren spielstylen zusammen. taktische spieler fanden es am unteren ende und wünschten sich mehr telegraphing. die agressiven spieler, die schaden wenig interessiert hat fanden es genug. Soldiers konnte man gut genug sehen da diese zum spieler laufen müssen um schaden zu machen aber das war nicht der fall bei den archers die von hinten ohne telegraphing angreiffen konnten. des weiteren gibt es zwei probleme die man trennen muss zum einen die wahrnehmung also das telegraphing an sich und dann auch nocht die reaktionsmöglichkeiten, wie man auf das telegraphing reagieren kann.

\paragraph{V5}
hier gibt es auch unterschiedliche meinungen, tp2 und 3 stimmen zu und tp3 erkennt muster in den wellen während tp2 die gegnerverteilungen in den räumen gefallen hat. dagegen sind aber tp4 und 5 die gefunden haben, dass zwar das dynamische verhalten der gegner sichtbar war aber die anfangsposition nicht unbedingt. manchmal sind auch archers und so vergessen gegangen. unerwartete Treffer von hinten gab es trotzdem und man kann die gegnerformation nicht sehen befor man in den raum tritt. die antworten sind des weiteren etwas schwer zu deuten, da die leitfrage 2 sachen gleichzeitig frägt, also die überschaubarkeit zu beginn und danach das unvorhersehbare dynamische verhalten.

\paragraph{V6: Der Raum als Druckmittel}
Dass TP3 in kleinen Räumen teils gleich aggressiv, teils hektischer spielte, liegt wahrscheinlich daran, dass sie wegen ihres aggressiven Spielstils ohnehin wenig Platz brauchte und den zusätzlichen Platz der Prospect Spaces nur in Momenten mit wenig Leben oder zu grosser Gefahr vermisste.

\paragraph{V7: Wellenstruktur mit Erholungsphasen}
Dass TP2 es schwer fand, Heilpausen zu finden, liegt daran, dass sie die Wellenstruktur des Spiels nicht erkannt hat und deshalb nicht davon profitierte. \textcolor{red}{[Belegt? Woran erkennt man, dass sie die Wellenstruktur nicht erkannt hat? Sonst vorsichtiger formulieren: \enquote{könnte daran liegen}.]}

\subsubsection{Gesamteindruck: Spielgefühl, Spielfluss und Schwierigkeit} \label{sec:Gesamteindruck}

\begin{itemize}[leftmargin=1.5cm]
    \item \textcolor{red}{Auswertung von Q9--Q11 und B14 zusammen mit den Ergebnissen pro Vorschlag. Das ist die Brücke zur Erkenntnisfrage in Kapitel~\ref{sec:Schlussfolgerungen}.}
\end{itemize}

\subsection{Verbesserungsvorschläge} \label{sec:Verbesserungen}

\begin{itemize}[leftmargin=1.5cm]
    \item \textcolor{red}{Konkrete Verbesserungen, jede mit Beleg aus den Daten (z.\,B. akustische Angriffsankündigung, Tutorial, Signal beim Öffnen der Tür, abbrechbare Combos, Block). Laut Vertrag müssen sie nicht mehr umgesetzt werden.}
\end{itemize}

\end{document}
